\section{Three-scale peeling and a near-\texorpdfstring{$11/21$}{11/21}
factorable wedge}
\label{sec:factorable-wedge}

The square-root normal form is best for locating the source-row wall.  To
use factorable-modulus distribution theorems, we return to the symmetric
TPC-15 packet with $U=V=\lfloor R^{1/2}\rfloor$ and open
\eqref{eq:betaV-intro}.  Writing $k=dj$ reveals the progression modulus
\begin{equation}
 q=\ell d.
 \label{eq:true-progression-modulus}
\end{equation}
Thus the relevant geometry has three factor scales $(L,K,D)$, with the
quotient length $J\asymp K/D$.

\subsection{The exact invisible wedge}

Define
\begin{equation}
 \gamma_{\ell;R,V}(k)=
 \sum_{\substack{d\mid k\\d\le V\\\ell d>R}}\mu(d).
 \label{eq:invisible-divisor-coefficient}
\end{equation}

\begin{proposition}[Hyperbolic Type-I peeling]
\label{prop:hyperbolic-peeling}
For every $A>0$,
\begin{align}
 \cH^{\rm HB}_{h,R;U,V}(X)
 ={}&\sum_{\substack{\ell>U\\k>V}}
 \Lambda(\ell)\gamma_{\ell;R,V}(k)
 r_R(\ell k+h)W(\ell k/X)
 \notag\\
 &+O_{A,h,W,\delta}\bigl(X(\log X)^{-A}\bigr).
 \label{eq:hyperbolically-peeled-packet}
\end{align}
In particular, if $U<\ell\le R$, only divisors $d>R/\ell$ survive; if
$\ell>R$, the coefficient is the full $\beta_V(k)$.
\end{proposition}

\begin{proof}
The difference between \eqref{eq:TPC15-hard-packet} and the main term in
\eqref{eq:hyperbolically-peeled-packet} is
\begin{equation}
 \sum_{\substack{\ell>U,\ d\le V\\\ell d\le R}}
 \Lambda(\ell)\mu(d)
 \sum_{j\ge1}r_R(\ell dj+h)W(\ell dj/X),
 \label{eq:visible-wedge}
\end{equation}
for all sufficiently large $X$.  Indeed, on this support
$\ell\le R$ and $\ell dj\asymp_WX$, so
$k=dj\gg_WX/R\gg V$; the original restriction $k>V$ is automatic.

Group \eqref{eq:visible-wedge} by $m=\ell d\le R$.  Its row coefficient
is
\[
 A_m=\sum_{\substack{\ell d=m\\\ell>U,\ d\le V}}
 \Lambda(\ell)\mu(d),
 \qquad |A_m|\le\sum_{\ell\mid m}\Lambda(\ell)=\log m.
\]
The full bounded Type-I row theorem of TPC-15, with one extra logarithmic
saving, makes \eqref{eq:visible-wedge}
$O_A(X(\log X)^{-A})$.
\end{proof}

The curve $\ell d=R$ is an exact calibration boundary for this
decomposition.  It is not an absolute boundary for all analytic methods;
factorable-modulus theorems can enter part of the exterior.

\subsection{An imported fixed-residue distribution theorem}

We use the following consequence of \citet[Theorem~1.1]{Maynard2025I}.
It is stated in a form adapted to smooth von Mangoldt sums.

\begin{theorem}[Maynard, fixed residue and factorable moduli]
\label{thm:Maynard-factorable}
Fix $a\ne0$, $\eta>0$, and $A>0$.  Suppose
\begin{equation}
 Q_1Q_2^2\le X^{1-\eta},\qquad
 Q_1^{12}Q_2^7\le X^{4-\eta},\qquad
 Q_1^{20}Q_2^{19}\le X^{10-\eta}.
 \label{eq:Maynard-three-inequalities}
\end{equation}
Let $G_{q_1,q_2}\in C_c^1((0,\infty))$ range over a family with a common
compact support and a common $C^1$ bound.  Then
\begin{align}
 &\sum_{\substack{q_1\le Q_1,\ q_2\le Q_2\\(q_1q_2,a)=1}}
 \left|
 \sum_{\substack{n\ge1\\n\equiv a\ ({\rm mod}\ q_1q_2)}}
 \Lambda(n)G_{q_1,q_2}((n-a)/X)
 -\frac X{\varphi(q_1q_2)}I_0(G_{q_1,q_2})
 \right|
 \notag\\
 &\hspace{16em}\ll_{a,\eta,A}X(\log X)^{-A}.
 \label{eq:Maynard-smoothed-consequence}
 \end{align}
\end{theorem}

\begin{proof}[Derivation of the stated smooth form]
Maynard's theorem is stated for
\[
 E(t;q,a)=\pi(t;q,a)-\frac{\pi(t)}{\varphi(q)}.
\]
Stieltjes partial summation, followed by the triangle inequality and
Fubini, bounds the sum of the prime-only weighted errors by a constant
times
\[
 (\log X)\sup_{t\asymp X}\sum_{q_1,q_2}|E(t;q_1q_2,a)|
 +\frac{\log X}{X}\int_{t\asymp X}
 \sum_{q_1,q_2}|E(t;q_1q_2,a)|\,dt.
\]
The common support and $C^1$ bounds make this estimate uniform even
though $G_{q_1,q_2}$ depends on the ordered modulus pair.  Apply
\citet[Theorem~1.1]{Maynard2025I} pointwise in $t\asymp X$, with its
parameter chosen smaller than $\eta/100$ and with two additional
logarithmic powers requested.  This gives
$O_A(X(\log X)^{-A})$.

The global prime number theorem, uniformly for the same bounded weight
family, gives
\[
 \sum_p(\log p)G_{q_1,q_2}((p-a)/X)
 =XI_0(G_{q_1,q_2})
 +O\left(X\exp(-c\sqrt{\log X})\right).
\]
Its aggregate error is negligible because
\[
 \sum_{q_1\le Q_1,q_2\le Q_2}
 \frac1{\varphi(q_1q_2)}
 \le
 \left(\sum_{q_1\le Q_1}\frac1{\varphi(q_1)}\right)
 \left(\sum_{q_2\le Q_2}\frac1{\varphi(q_2)}\right)
 \ll(\log X)^C.
\]
Finally, restore target prime powers.  If $p^\nu\asymp X$, $\nu\ge2$,
and $p^\nu\equiv a\pmod{q_1q_2}$, then $q_1q_2\mid p^\nu-a$.
Consequently their total over the ordered pairs is
\[
 \ll\sum_{\substack{p^\nu\asymp X\\\nu\ge2}}
 (\log p)\tau_3(|p^\nu-a|)
 \ll_\eps X^{1/2+\eps}.
\]
This proves \eqref{eq:Maynard-smoothed-consequence}.  The absolute values
in Maynard's theorem are what permit dyadic restrictions and bounded
outer coefficients.  No new distribution estimate is proved here.
\end{proof}

Let $\Phi_1,\Phi_2,\Phi_3\in C_c^\infty([1,2])$.  Define the smooth
$D$-resolved hard subpacket
\begin{align}
 \cH_{h,R}(L,K,D)
 ={}&\sum_{\ell,d,j\ge1}
 \Lambda(\ell)\mu(d)
 \Phi_1(\ell/L)\Phi_2(d/D)\Phi_3(dj/K)
 \notag\\
 &\hspace{4.2em}\times
 r_R(\ell dj+h)W(\ell dj/X).
 \label{eq:three-scale-subpacket}
\end{align}

\begin{theorem}[A factorable exterior wedge]
\label{thm:factorable-exterior-wedge}
Fix $h\ne0$, $0<\delta<1/2$, and $\eta>0$.  Suppose the scales satisfy
\begin{align}
 &LK\asymp_WX,
 \qquad 2D<R<L/2,
 \qquad 2D\le V,
 \qquad K>2V,
 \qquad L>2U,
 \label{eq:wedge-geometric-conditions}\\
 &DL^2\le X^{1-\eta},
 \qquad D^{12}L^7\le X^{4-\eta},
 \qquad D^{20}L^{19}\le X^{10-\eta},
 \label{eq:wedge-Maynard-conditions}\\
 &LDR\le X^{1-\eta}.
 \label{eq:wedge-leakage-condition}
\end{align}
Then for every $A>0$,
\begin{equation}
 \cH_{h,R}(L,K,D)
 \ll_{A,h,W,\Phi_1,\Phi_2,\Phi_3,\delta,\eta}
 X(\log X)^{-A}.
 \label{eq:factorable-wedge-bound}
\end{equation}
\end{theorem}

\begin{proof}
The contribution of source prime powers $\ell=p^a$, $a\ge2$, is removed
first.  The divisor expansion gives $|r_R(n)|\ll_\eps X^\eps$ on the
support in question.  Hence their entire subpacket is
\begin{equation}
 \ll_\eps
 L^{1/2}D\frac KD X^\eps
 \ll X L^{-1/2+\eps},
 \label{eq:source-prime-power-wedge}
\end{equation}
because $LK\asymp X$.  This has a fixed power saving since $L>2U$.
We may therefore assume below that $\ell$ is prime.  In particular,
$\ell>L>2R>|h|$ for all sufficiently large $X$.

Fix $\ell\asymp L$ and $d\asymp D$, and set $q=\ell d$.  The inner sum
is a smooth progression row of modulus $q$, with a family of weights
whose supports and $C^1$ norms are uniform because $LK\asymp X$.
Equation \eqref{eq:row-drift-error-leakage} decomposes it into prime error,
cutoff drift, and model leakage.

On rows with $(\ell d,h)=1$, apply
\cref{thm:Maynard-factorable} with $Q_1=2D,Q_2=2L$.  The powers of $2$ are
absorbed by the strict $X^{-\eta}$ margins.  The theorem sums absolute
errors, so the restrictions $d\asymp D$, $\ell\asymp L$ and the weights
$|\mu(d)|$, $\Lambda(\ell)\ll\log X$ cost only a change in the requested
logarithmic saving.  This bounds the total prime-progression error by
$O_A(X(\log X)^{-A})$.

If $(\ell d,h)>1$, the common prime lies in $d$, because the remaining
$\ell$ is prime and exceeds $|h|$.  If $p\mid(d,h)$, then
$\ell dj+h$ is divisible by $p$; a nonzero von Mangoldt value therefore
forces $\ell dj+h=p^\nu$.  For each of the finitely many $p\mid h$, only
$O_{h,W}(1)$ exponents lie on the $X$ scale, and the number of ordered
factorizations $\ell d\mid p^\nu-h$ is $O_\eps(X^\eps)$.  Thus the total
prime part of the nonreduced rows is negligible.  Their complete model
mean is zero: all divisors of $d$ lie below $R$, and their local Euler
factor contains $1-1$ at a prime dividing $(d,h)$.

The finite-row model leakage is
$O(R(\log(2R))^2)$ per pair $(\ell,d)$.  After the outer logarithmic
weight and the $O(LD)$ pairs, its total is
\begin{equation}
 \ll LDR(\log X)^3,
 \label{eq:aggregate-model-leakage-wedge}
\end{equation}
which has a fixed power saving by \eqref{eq:wedge-leakage-condition}.

It remains to treat the drift.  Since $d< R<\ell$ and $\mu(d)\ne0$, every
divisor of $\ell d$ below $R$ is a divisor of $d$.  Therefore
\begin{equation}
 \Delta_{h,R}(\ell d)=
 \begin{cases}
 \displaystyle\frac{d}{\varphi(d)(\ell-1)},&(d,h)=1,\\[0.8em]
 0,&(d,h)>1.
 \end{cases}
 \label{eq:simplified-wedge-drift}
\end{equation}
Indeed, the model mean is $d/\varphi(d)$ and the prime mean is
$\ell d/(\varphi(\ell)\varphi(d))$ in the first case; both are zero in
the second.  The continuous drift contribution is consequently
\begin{align}
 &\ll X
 \sum_{\ell\asymp L}
 \frac{\Lambda(\ell)}{\ell(\ell-1)}
 \sum_{d\asymp D}\frac{\mu(d)^2}{\varphi(d)}
 \notag\\
 &\ll \frac X L(\log(2D))^2.
 \label{eq:wedge-drift-total}
\end{align}
This also has a fixed power saving.  Combining the three pieces proves
\eqref{eq:factorable-wedge-bound}.
\end{proof}

\begin{corollary}[A packet near the \texorpdfstring{$11/21$}{11/21}
factorable boundary]
\label{cor:eleven-twenty-one-packet}
Fix $0<\sigma<1/100$ and
\begin{equation}
 \frac1{42}+3\sigma<\delta<\frac14.
 \label{eq:delta-sigma-range}
\end{equation}
Take
\begin{equation}
 D=X^{1/21-\sigma},
 \qquad
 L=X^{10/21-\sigma},
 \qquad
 K=X/L=X^{11/21+\sigma}.
 \label{eq:explicit-wedge-scales}
\end{equation}
Then, taking $\eta=2\sigma$ in
\cref{thm:factorable-exterior-wedge}, the packet
\eqref{eq:three-scale-subpacket} is
$O_A(X(\log X)^{-A})$ for every $A>0$.
\end{corollary}

\begin{proof}
The three Maynard monomials have exponents
\[
 1-3\sigma,
 \qquad \frac{82}{21}-19\sigma<4,
 \qquad 10-39\sigma.
\]
Moreover,
\[
 LDR=X^{1+1/42-2\sigma-\delta}
 \le X^{1-5\sigma}
\]
by \eqref{eq:delta-sigma-range}.  The same range gives $R<L$; the upper
bound $\delta<1/4$ ensures $2D<V$ for large $X$.  All remaining support
conditions in \cref{thm:factorable-exterior-wedge} are immediate.
\end{proof}

Both $L$ and $K$ exceed $R$ in this corollary, so the packet lies inside
the central $(L,K)$ corridor of the original TPC-15 representation.  The
result is nevertheless only one $D$-resolved slice.  Summing over all
$D$, or removing the Maynard factorization inequalities, is not justified
by the theorem.
